\documentclass[10pt,a4paper]{amsart}
\usepackage[margin=19mm]{geometry}
\usepackage{amsmath,amssymb,mathtools}
\usepackage[scr=rsfs,cal=euler]{mathalfa}
\usepackage{xfrac}
\usepackage[unicode,hidelinks,bookmarks=false]{hyperref}
\newcommand{\ie}{i.e.\ }
\newcommand{\cf}{cf.\ }
\newcommand{\resp}{respectively\ }
\newcommand{\Subsection}{\subsection}
\providecommand{\nobreakdash}{\nobreak\hbox{-}}
\setlength{\emergencystretch}{2em}
\multlinegap0pt
\usepackage{bm}
\usepackage{bbm}
\usepackage{dsfont}
\usepackage{xspace}
\usepackage{cancel}
\usepackage{mathtools}
\usepackage{amsthm}
\makeatletter
\@ifundefined{thm}{\newtheorem{thm}{Theorem}}{}
\@ifundefined{assu}{\newtheorem{assu}[thm]{Assumption}}{}
\@ifundefined{lem}{\newtheorem{lem}[thm]{Lemma}}{}
\@ifundefined{prop}{\newtheorem{prop}[thm]{Proposition}}{}
\makeatother
\newcommand{\natu}{\mathbb N}
\newcommand{\nat}{\mathbb N_0}
\newcommand{\rzm}{\mathbb R^{m}}
\newcommand{\rzn}{\mathbb R^{n}}
\newcommand{\rz}{\mathbb R}
\title[The maximum principle for discrete-time control systems and ]{The maximum principle for discrete-time control systems and applications to dynamic games}
\author{Alberto Dom\'inguez Corella, On\'esimo Hern\'andez-Lerma}
\date{Source: arXiv:2601.11395v1 (CC BY 4.0)}
\begin{document}
\maketitle

\noindent\emph{Source and licence.} The maximum principle for discrete-time control systems and applications to dynamic games. Version arXiv:2601.11395v1 (CC BY 4.0), distributed under the Creative Commons Attribution 4.0 International licence, as recorded in the arXiv licence metadata for this exact version and shown on the abstract page https://arxiv.org/abs/2601.11395v1. Full rights notice: licenses/arxiv-2601.11395-CC-BY-4.0.txt.\par\smallskip

\noindent\textit{Editorial scope.} Counted unit: the Thm whose source label is \texttt{MP} in the article (source0.tex lines 320--340, source label \texttt{MP}), delivered together with the article's own complete original proof of it (source0.tex lines 341--371, one \texttt{proof} environment of 31 lines). No mathematics is written, repaired or re-derived: statement and proof are the article's own text line for line. Required context retained uncounted: dynamic (lines 168--170); tuple (lines 192--194); PI (lines 217--219); GD (lines 258--263); AMP (lines 281--285); DV (lines 304--317). Every definition, display and label of those lines is retained unchanged. Omitted from this package and not redistributed: 1--167, 171--191, 195--216, 220--257, 264--280, 286--303, 318--319, 372--1545. Nothing inside a retained excerpt is omitted or altered: every retained body is delivered exactly as the article's lines are written, so no omission receipt of any kind accompanies this package. Citations: the retained excerpts carry no citation command at all, so no bibliography entry of the article is transcribed and no reference is redistributed. Reference adaptation: none was needed and none was made; every reference inside the delivered unit resolves inside it, so no unresolved reference or citation is produced. Printed slips kept literal: no printed word, symbol, hypothesis or number of the counted statement or of its proof was altered. Layout and class: the article's own class is \texttt{elsarticle} with packages including amsmath, mathrsfs, amssymb, amsfonts, amsthm, newtxmath; the wrapper is the lane's pinned \texttt{amsart} editorial wrapper at 10pt on A4, which restates the theorem-like environments the retained text needs and adds the article's own macro definitions that the retained text needs (\texttt{\textbackslash nat}, \texttt{\textbackslash natu}, \texttt{\textbackslash rz}, \texttt{\textbackslash rzm}, \texttt{\textbackslash rzn}), with extra wrapper packages (bm, bbm, dsfont, xspace, cancel, mathtools). The delivered numbering is the unit's own. Assets: no retained span contains a figure environment, a graphics-inclusion command, a PDF-page inclusion, a table or a TikZ picture; no image file is copied into this package, the package declares no asset dependency and no image file is redistributed with it. Licence: version arXiv:2601.11395v1 of this article is distributed under the Creative Commons Attribution 4.0 International licence, as recorded in the arXiv licence metadata for this exact version and shown on the abstract page https://arxiv.org/abs/2601.11395v1. A full rights notice is delivered at licenses/arxiv-2601.11395-CC-BY-4.0.txt. Characters: every retained excerpt is pure ASCII, so the delivered engine, XeLaTeX with its default Latin Modern text font, reports no missing character.
	\begin{align}\label{dynamic}
	x_{t+1}=f_t(x_t,u_t) \hspace{.5cm}\forall_{t\in\nat},
	\end{align}

	\begin{align}\label{tuple}
	(\Psi, \left\lbrace f_t\right\rbrace, \left\lbrace g_t\right\rbrace),
	\end{align} 

	\begin{align}\label{PI}
	v(\psi)=\sum_{t=0}^\infty g_t(x_t^\psi,u_t).
	\end{align}

	\begin{prop}\label{GD}
		Let $\mathcal{X}$ be a linear space and $\mathcal{V}$ a subset of $\mathcal X$. Let $p\in\mathcal V$ be an internal point in the direction $q\in\mathcal X$ and $h:\mathcal V\to\rz$ a given function. If the Gateaux differential of $h$ at $p$ in the  direction $q$ exists and $h$ has a maximum at $p$, then 
		\begin{align*}
		\delta_h(p;q)=0.
		\end{align*}
	\end{prop}

	\begin{assu}\label{AMP}
		Let  $\hat \psi=(\hat u_0,\hat u_1,\dots) \in \Psi$. For each  $\tau\in\nat$, define the sequence of functions $\rho^\tau_t: U_\tau\to\rzn$ as $$\rho_t^\tau(u):=\frac{\partial g_t}{\partial x}( x_t^{{\hat\psi^\tau(u)}}, \hat u_t)\prod_{s=\tau+1}^{t-1}\frac{\partial f_s}{\partial x}(x_s^{{\hat \psi^\tau(u)}}, \hat u_s),$$ 
		where $\hat\psi^\tau(u)=(\hat u_0,\dots,\hat u_{\tau-1},u,\hat u_{\tau+1},\dots)$. 
		We suppose that, for each $\tau\in\nat$, there exists an open neighborhood $\mathcal O_\tau\subset U_\tau$ of $\hat u_\tau$ such that the series $\displaystyle\sum_{t=\tau+1}^{\infty}\rho^\tau_t$ converges uniformly on $\mathcal O_\tau$. 
	\end{assu}

	\begin{lem}\label{DV}
		Let  $\hat \psi=\left\lbrace \hat u_t\right\rbrace \in \Psi$ be a plan for which Assumption \ref{AMP} holds. Let $y\in\rzm$ and $\tau\in\nat$. Then $\hat\psi$ is an internal point of $\Psi$ in the direction $\psi^{\tau,y}\in\Lambda$, where $\psi^{\tau,y}$ is defined as 
		\begin{align*}
		\psi^{\tau,y}_t:=\left\{ \begin{array}{lcc}
		\displaystyle y &   if  & t=\tau \\
		\\ 0 &  if & t\neq\tau, \\
		\end{array}
		\right.
		\end{align*}
		for all $t\in\nat$. Moreover, if $v$ is the function in (\ref{PI}), its G\^ ateaux differential at $\hat\psi$ in the direction $\psi^{\tau,y}$ exists and is given by 
		\begin{align*}
		\delta_v(\hat\psi;\psi^{\tau,y})= \left( \sum_{t=\tau+1}^\infty\frac{\partial g_t}{\partial x}(x_t^{\hat{\psi}},\hat u_t)\prod_{s=\tau+1}^{t-1}\frac{\partial f_s}{\partial x}(x_s^{\hat{\psi}},\hat u_s)\right)\frac{\partial f_\tau}{\partial y}(x_\tau^{\hat{\psi}},\hat u_\tau)y^*+\frac{\partial g_\tau}{\partial y}(x_\tau^{\hat{\psi}},\hat u_\tau)y^*.
		\end{align*} 
	\end{lem} 

	\begin{thm}\label{MP}
		Let $\hat \psi=\left\lbrace\hat u_t \right\rbrace \in\Psi$ be a plan for which Assumption \ref{AMP} holds. If $\hat\psi$ is an optimal plan of the OCP (\ref{dynamic})-(\ref{tuple}), then there exists a sequence $\left\lbrace \lambda_t\right\rbrace_{t=1}^\infty$ in $\rzn$ such that
		\begin{enumerate}
			\item[(a)] For all $t\in\natu$,
			\begin{align}\label{MPX}
			\frac{\partial g_t}{\partial x}(x_t^{\hat\psi},\hat u_t)+\lambda_{t+1}\frac{\partial f_t}{\partial x}(x_t^{\hat\psi},\hat u_t)=\lambda_t,
			\end{align}
			\item[(b)] For all $t\in\nat$,
			\begin{align}\label{MPY}
			\frac{\partial g_t}{\partial y}(x_t^{\hat\psi},\hat u_t)+\lambda_{t+1}\frac{\partial f_t}{\partial y}(x_t^{\hat\psi},\hat u_t)=0,
			\end{align}
			\item[(c)] For every $h\in\natu$, we have the  transversality condition (TC)
			\begin{align}\label{TC}
			\lim_{t\to\infty}\lambda_t \prod_{s=h}^{t-1}\frac{\partial f_s}{\partial x}(x_s^{\hat\psi},\hat u_s)=0.
			\end{align}  
		\end{enumerate}
		Moreover, each $\lambda_t$ is given by 
		\begin{align}\label{multipliers}
		\lambda_t=\sum_{k=t}^\infty\frac{\partial g_k}{\partial x}(x_k^{\hat{\psi}},\hat u_k)\prod_{s=t}^{k-1}\frac{\partial f_s}{\partial x}(x_s^{\hat{\psi}},\hat u_s).
		\end{align}
	\end{thm}

	\begin{proof}
		
		\begin{enumerate}
			\item[(a)] Pick an arbitrary $\tau\in\natu$. Then
			\begin{align*}
			\lambda_\tau&:=\sum_{t=\tau}^\infty\frac{\partial g_t}{\partial x}(x_t^{\hat{\psi}},\hat u_t)\prod_{s=\tau}^{t-1}\frac{\partial f_s}{\partial x}(x_s^{\hat{\psi}},\hat u_s)\\
			&=\frac{\partial g_\tau}{\partial x}(x_\tau^{\hat{\psi}},\hat u_\tau)+\sum_{t=\tau+1}^\infty\frac{\partial g_t}{\partial x}(x_t^{\hat{\psi}},\hat u_t)\prod_{s=\tau}^{t-1}\frac{\partial f_s}{\partial x}(x_s^{\hat{\psi}},\hat u_s)\\
			&=\frac{\partial g_\tau}{\partial x}(x_\tau^{\hat{\psi}},\hat u_\tau)+\left( \sum_{t=\tau+1}^\infty\frac{\partial g_t}{\partial x}(x_t^{\hat{\psi}},\hat u_t)\prod_{s=\tau+1}^{t-1}\frac{\partial f_s}{\partial x}(x_s^{\hat{\psi}},\hat u_s)\right)\frac{\partial f_\tau}{\partial x}(x_\tau^{\hat{\psi}},\hat u_\tau)\\
			&=\frac{\partial g_\tau}{\partial x}(x_\tau^{\hat{\psi}},\hat u_\tau)+\lambda_{\tau+1}\frac{\partial f_\tau}{\partial x}(x_\tau^{\hat{\psi}},\hat u_\tau).
			\end{align*}
			\item[(b)] Fix $\tau\in\nat$ and $y\in\rzm$ arbitrary. By Lemma \ref{DV} and Proposition \ref{GD},
			\begin{align*}
			\left( \sum_{t=\tau+1}^\infty\frac{\partial g_t}{\partial x}(x_t^{\hat{\psi}},\hat u_t)\prod_{s=\tau+1}^{t-1}\frac{\partial f_s}{\partial x}(x_s^{\hat{\psi}},\hat u_s)\right)\frac{\partial f_\tau}{\partial y}(x_\tau^{\hat{\psi}},\hat u_\tau)y^*+\frac{\partial g_\tau}{\partial y}(x_\tau^{\hat{\psi}},\hat u_\tau)y^*=0,
			\end{align*}
			that is 
			\begin{align*}
			\left[\frac{\partial g_\tau}{\partial y}(x_\tau^{\hat{\psi}},\hat u_\tau)y^*+\lambda_{\tau+1}\frac{\partial f_\tau}{\partial y}(x_\tau^{\hat{\psi}},\hat u_\tau)y^*\right]=0.
			\end{align*}
			Since this holds for any $y\in\rzm$, $(b)$ follows.
			\item[(c)] By Assumption \ref{AMP}, the series
			\begin{align*}
			\sum_{k=h}^\infty\frac{\partial g_k}{\partial x}(x_k^{\hat{\psi}},\hat u_k)\prod_{s=h}^{k-1}\frac{\partial f_s}{\partial x}(x_s^{\hat{\psi}},\hat u_s),
			\end{align*}
			converges for any $h\in\natu$ and
			\begin{align*}
			\lambda_t \prod_{s=h}^{t-1}\frac{\partial f_s}{\partial x}(x_s^{\hat\psi},\hat u_s)&=\left( \sum_{k=t}^\infty\frac{\partial g_k}{\partial x}(x_k^{\hat{\psi}},\hat u_k)\prod_{s=t}^{k-1}\frac{\partial f_s}{\partial x}(x_s^{\hat{\psi}},\hat u_s)\right)  \prod_{s=h}^{t-1}\frac{\partial f_s}{\partial x}(x_s^{\hat\psi},\hat u_s)\\
			&=\sum_{k=t}^\infty\frac{\partial g_k}{\partial x}(x_k^{\hat{\psi}},\hat u_k)\prod_{s=h}^{k-1}\frac{\partial f_s}{\partial x}(x_s^{\hat{\psi}},\hat u_s).
			\end{align*}
			Letting $t$ tend to infinity, $(c)$ follows.
		\end{enumerate}
	\end{proof}

\end{document}
